\section{Related Work}
\label{sec:related}

\textbf{Docking with moving stations.} Docking guidance is staged by sensing modality, with acoustic or electromagnetic homing bringing the vehicle within camera range and vision closing the final metres \cite{yazdani2020survey,maire2009vision}. In nearly all of it the dock is a fixed pose and the disturbance a steady current on the vehicle. Trsli\'c et al.\ docked a work-class ROV to a tether management system heaving 1.1~m peak to peak at 8.5~s by servoing on the latest camera pose, holding the last pose through dropouts and setting the depth setpoint to the cage's mean depth; they report that the vehicle's inertia made compensation infeasible and contact inevitable \cite{trslic2020vision}. Their neuro-fuzzy predictor forecast the cage depth 2.5~s ahead, from a pressure sensor on the cage, for heave only, offline \cite{trslic2020neuro}. A seabed-moored suspended dock under steady current swings by more than a metre and yaws without settling; its designers address this with structure and a funnel, not estimation \cite{guo2024dynamic}. At a shallow test site wave-induced forces dislodged a resident vehicle from its dock several metres below the surface; its 90~m operational site is free of them \cite{gunzel2026docking}. Trsli\'c et al.\ did predict the cage's wave-frequency motion, from a sensor on the cage; none of these estimates it through the vehicle's camera or feeds it forward.

\textbf{Estimating a moving target from a moving observer.} Target tracking supplies the kinematic hierarchy: a constant-position model absorbs target motion as process noise, a constant-velocity model adds a velocity state driven by white acceleration, and the acceleration density sets the trade between responsiveness and noise amplification \cite{barshalom2001estimation}. Every member assumes white measurement errors; correlated errors need state augmentation or differencing, and no process-noise setting substitutes for it \cite{bryson1968correlated}. When the target is measured relative to the observer, the natural formulation is a relative state with the observer's velocity as a known input, as in spacecraft relative navigation \cite{kim2007relative}, so the observer's velocity error enters the process model. Wave-driven platform motion has been estimated as a second-order oscillator for ships \cite{triantafyllou1983ship} and as an internal model for landing on an oscillating deck \cite{marconi2002landing}; the constant-velocity model used here is deliberately model-free, because a real dock's motion is not a single sinusoid. In the docking literature the filter is a noise suppressor for a static dock \cite{gunzel2026docking,ni2025vision}, or a disturbance observer on the vehicle with a constant-input random-walk state \cite{shi2023integrated}; the case of a velocity state on the target, fed forward, is not examined.

\textbf{Visual servoing and latency.} Position-based visual servoing regulates a reconstructed relative pose, which costs nothing extra when the marker pipeline already produces one \cite{chaumette2006visual,garrido2014automatic}. Tracking a moving target adds a feedforward of the estimated target velocity to the servo law \cite{chaumette2007visual}, and the closed-loop behaviour then depends on the dynamics between the command and the vehicle's response \cite{corke1996dynamic}. On the vehicle side the command channel of the production autopilot is effort-like rather than a velocity setpoint, so the feedforward gain must be matched to an identified plant gain, a point that the failure analysis of Section~\ref{sec:failure} turns on. The remedy, a cascade with an inner velocity loop on the vehicle's measured velocity and the target velocity added to its reference, is the standard structure of marine motion control \cite{fossen2011handbook} and of velocity feedforward in visual servoing \cite{corke1996dynamic}; the contribution here is the evidence of what the open-loop law does on an effort-like channel and of what the inner loop buys under navigation error.
